\section{Purpose, scope, and logical status}
\label{sec:scope}

This document converts the project discussion into a mathematical research
note suitable for incorporation into the repository.  It is deliberately not
a verbatim transcript.  The discussion contained exploratory claims, notation
collisions, and one important error concerning the SOR relaxation range.  The
present note performs four operations:
\begin{enumerate}[label=(\roman*),leftmargin=2.2em]
  \item it restates the two source frameworks in a common normalization;
  \item it separates empirical motivation from theorem statements;
  \item it proves the parts of the hybrid analysis that are currently closed;
  \item it records the exact remaining locality lemmas needed for universal
        accelerated-work theorems.
\end{enumerate}

The first part studies the unregularized symmetric PageRank quadratic used by
\citet{huang2025accelerated}.  The second part treats the
$\ell_1$-regularized PageRank objective of
\citet{fountoulakis2026complexity} as a separate composite optimization
problem.  We do not silently identify the two objectives: RPPR support and KKT
results are used only where their assumptions apply, while their confinement
techniques also motivate explicit conditional results for unregularized AESP.

\begin{table}[t]
\centering
\caption{Status of the principal statements in this note.}
\label{tab:status-overview}
\begin{tabularx}{\textwidth}{@{}>{\raggedright\arraybackslash}p{0.39\textwidth}>{\raggedright\arraybackslash}p{0.16\textwidth}X@{}}
\toprule
Statement & Status & Qualification \\
\midrule
Locally evolving-set and AESP source guarantees
  & \statussource & Restated from
    \citet{zhou2024iterative,huang2025accelerated}. \\
AESP is especially effective in the early stage
  & \statusobservation & Supported by Figure~4 of
    \citet{huang2025accelerated}; not a universal theorem. \\
Finite convergence of AESP followed by local SOR
  & \statusproved & Holds from every finite handoff for fixed
    $0<\omega<2$. \\
$\bigO(1/(\sqrt\alpha\eps))$ SOR tail after an objective-gap handoff
  & \statusproved & Uses $\omega=1$ and the handoff
    $f(x_J)-f^\star\leq\alpha^{3/2}\eps/(1+\alpha)$. \\
$\bigO(1/(\sqrt\alpha\eps))$ SOR tail after a weighted-$\ell_1$ handoff
  & \statusproved & Uses $\omega=1$ and
    $\norm{D^{1/2}\nabla f(x_J)}_1=\bigO(\alpha^{3/2})$. \\
Master handoff and end-to-end trajectory bounds
  & \statusproved & Depend on the Phase-I work and, for the AESP trajectory
    form, on the realized early locality factor $\Lambda_J$. \\
Black-box $\bigO(R/(\alpha^{3/4}\eps))$ AESP--LOCSOR rate
  & \textbf{Parameterized} & Closed for inner-plus-tail work with realized
    AESP ratio $R$; outer initialization must be charged separately. \\
Graph-uniform $\softO(1/(\sqrt\alpha\eps))$ total work
  & \statusconditional & Follows from
    $\Lambda_J=\bigO(1/\eps)$; not yet proved for every graph. \\
Composite Catalyst followed by full ISTA for RPPR
  & \statusproved & The proximal map is a weighted
    $(1-\alpha)$ contraction, so every finite handoff converges. \\
Graph-uniform $\softO(1/(\rho\sqrt\alpha))$ RPPR work
  & \statusopen & Requires a local-work theorem from an arbitrary
    accelerated warm start. \\
Plain optimal local SOR is weighted-$\ell_1$ monotone
  & \textbf{False} & The universal monotone range is
    $0<\omega<1+\alpha$, not all of $(0,2)$. \\
Plain local Chebyshev gives an unconditional tail theorem
  & \statusopen & The source theorem retains a
    localization-forcing condition. \\
\bottomrule
\end{tabularx}
\end{table}

\section{Problem, normalization, certificate, and work model}
\label{sec:problem}

\subsection{Symmetric PageRank quadratic}

Let $G=(V,E)$ be a finite, connected, undirected simple graph without isolated
vertices.  Let $A\in\R^{n\times n}$ be its adjacency matrix,
$D=\diag(d_1,\ldots,d_n)$ its degree matrix, and
\[
  W:=D^{-1/2}AD^{-1/2}.
\]
Fix a source vertex $s\in V$ and a PageRank parameter
$\alpha\in(0,1)$.  We use the AESP normalization
\begin{equation}
  Q:=\frac{1+\alpha}{2}I-\frac{1-\alpha}{2}W
    =\alpha I+\frac{1-\alpha}{2}(I-W),
  \qquad
  b:=\alpha D^{-1/2}e_s.
  \label{eq:Q-b}
\end{equation}
The objective is
\begin{equation}
  f(x):=\frac12 x^{\trans}Qx-b^{\trans}x,
  \qquad
  g(x):=\nabla f(x)=Qx-b.
  \label{eq:f-g}
\end{equation}
Since the spectrum of $I-W$ is contained in $[0,2]$,
\begin{equation}
  \alpha I\preceq Q\preceq I.
  \label{eq:spectrum-Q}
\end{equation}
Thus $f$ is $\alpha$-strongly convex and $1$-smooth, and its unique minimizer
is $x^\star=Q^{-1}b$.  The corresponding personalized PageRank vector is
\begin{equation}
  \pi:=D^{1/2}x^\star.
  \label{eq:pi-from-x}
\end{equation}

\subsection{Coordinatewise accuracy certificate}

The AESP paper proves the following certificate in this normalization.

\begin{lemma}[Degree-normalized gradient certificate]
\label{lem:certificate}
For every $x\in\R^n$,
\begin{equation}
  \norm{D^{-1/2}(x-x^\star)}_\infty
  \leq
  \frac1\alpha\norm{D^{-1/2}g(x)}_\infty.
  \label{eq:certificate}
\end{equation}
Consequently,
\begin{equation}
  \norm{D^{-1/2}g(x)}_\infty\leq\alpha\eps
  \quad\Longrightarrow\quad
  \norm{D^{-1}(D^{1/2}x-\pi)}_\infty\leq\eps.
  \label{eq:certificate-stop}
\end{equation}
\end{lemma}

\begin{proof}
Let
$\Pi_\alpha:=\alpha D^{1/2}Q^{-1}D^{-1/2}$ be the PageRank matrix.  The source
analysis gives $\norm{D^{-1}\Pi_\alpha D}_\infty=1$.  Since
\[
  x-x^\star=Q^{-1}g(x),
  \qquad
  D^{-1/2}Q^{-1}D^{1/2}
  =\frac1\alpha D^{-1}\Pi_\alpha D,
\]
we obtain
\[
  \norm{D^{-1/2}(x-x^\star)}_\infty
  \leq
  \norm{D^{-1/2}Q^{-1}D^{1/2}}_\infty
  \norm{D^{-1/2}g(x)}_\infty
  =\frac1\alpha\norm{D^{-1/2}g(x)}_\infty.
\]
The second claim follows from $D^{-1}(D^{1/2}x-\pi)=D^{-1/2}(x-x^\star)$.
\end{proof}

It is convenient to define the degree-shifted gradient
\begin{equation}
  q(x):=D^{1/2}g(x).
  \label{eq:q-def}
\end{equation}
Then a coordinate is active at final tolerance $\eps$ precisely when
\begin{equation}
  \abs{g_u(x)}\geq\alpha\eps\sqrt{d_u}
  \quad\Longleftrightarrow\quad
  \abs{q_u(x)}\geq\alpha\eps d_u.
  \label{eq:final-active}
\end{equation}

\subsection{Dictionary between the 2024 and 2025 normalizations}

The locally evolving-set paper uses
\begin{equation}
  \widehat Q:=I-\chi W,
  \qquad
  \widehat b:=\frac{2\alpha}{1+\alpha}D^{-1/2}e_s,
  \qquad
  \chi:=\frac{1-\alpha}{1+\alpha}.
  \label{eq:work1-normalization}
\end{equation}
The two systems have the same solution because
\begin{equation}
  Q=\frac{1+\alpha}{2}\widehat Q,
  \qquad
  b=\frac{1+\alpha}{2}\widehat b.
  \label{eq:normalization-scale}
\end{equation}
If
$\widehat r:=\widehat b-\widehat Qx$, then
\begin{equation}
  g(x)=-\frac{1+\alpha}{2}\widehat r.
  \label{eq:gradient-residual-dictionary}
\end{equation}
Hence the Work-I residual certificate
$\norm{D^{-1/2}\widehat r}_\infty\leq 2\alpha\eps/(1+\alpha)$
is exactly the AESP certificate in \cref{eq:certificate-stop}.

\subsection{Degree-weighted local work}

Scanning the adjacency list of a vertex $u$ costs $d_u$ units.  A sequence of
coordinate updates $(u_j)_{j=0}^{K-1}$ therefore has work
\begin{equation}
  W:=\sum_{j=0}^{K-1}d_{u_j}.
  \label{eq:coordinate-work}
\end{equation}
For a parallel evolving-set iteration with active set $S_k$, the corresponding
work is $\vol(S_k):=\sum_{u\in S_k}d_u$.  This is the common unit used by
\APPR, \LOCSOR, the AESP inner methods, and the present hybrid analysis.
